\begin{figure}[htbp]
\centering
\begin{tikzpicture}[>=Latex,node distance=8mm]
\node[draw,rounded corners,align=center] (f) {local equation\\\(f=f_m+f_{m+1}+\cdots\)};
\node[draw,rounded corners,align=center,right=of f] (sub) {substitute\\\(y=xu\)};
\node[draw,rounded corners,align=center,right=of sub] (fac) {\(x^m\)\\exceptional factor};
\node[draw,rounded corners,align=center,right=of fac] (tan) {on \(E\):\\\(f_m(1,u)=0\)};
\draw[->] (f)--(sub);
\draw[->] (sub)--(fac);
\draw[->] (fac)--(tan);
\end{tikzpicture}
\caption{Multiplicity becomes exceptional multiplicity, while the tangent cone controls \(\widetilde C\cap E\).}
\label{fig:vi45-tangent-cone}
\end{figure}
